% Generated from locale/tr/content/lambda-calculus/syntax/free-variables.tex; do not hand-edit.


\olfileid[tr]{lam}{syn}{fv}
\olsection{Serbest Değişkenler}

Lambda hesabı fonksiyonları ele alır; fonksiyonlar lambda soyutlamasıyla
ortaya çıkar. Sezgisel olarak $\lambd[x][M]$, fonksiyonun argümanı $x$'e
atandığında değerleri $M$ tarafından verilen fonksiyondur. Ancak $M$ içindeki
her $x$ geçişi bu soyutlamayla ilgili değildir: $M$ başka bir
$\lambd[x][N]$ soyutlaması içeriyorsa $N$ içindeki $x$ geçişleri
$\lambd[x][N]$ ile ilgilidir, $\lambd[x][M]$ ile değil. Dolayısıyla
$\lambd[x][M]$ içindeki bir $\lambd[x]$ lambda soyutlaması, $M$ içindeki
başka bir lambda soyutlaması tarafından zaten bağlanmamış $x$ geçişlerini,
yani $M$ içindeki \emph{serbest} $x$ geçişlerini \emph{bağlar}.

\begin{defn}[Kapsam]
$\lambd[x][M]$ bir $N$ teriminin içinde geçiyorsa, $M$'nin karşılık gelen
geçişi $\lambd[x]$'in \emph{kapsamıdır}.
\end{defn}


\begin{defn}[Serbest ve bağlı geçiş]
  Bir $M$ terimindeki $x$ değişkeninin bir geçişi, herhangi bir
  $\lambd[x]$'in kapsamında değilse \emph{serbest}, aksi durumda
  \emph{bağlıdır}. $\lambd[x][M]$ içindeki bir $x$ geçişi, ancak ve ancak
  $M$ içindeki karşılık gelen $x$ geçişi serbestse baştaki $\lambd[x]$
  tarafından bağlanır.
\end{defn}

\begin{ex}
  $\lambd[x][x y]$ içinde hem $x$ hem de $y$, $\lambd[x]$'in kapsamındadır;
  dolayısıyla $x$, $\lambd[x]$ tarafından bağlanır. $y$ herhangi bir
  $\lambd[y]$'nin kapsamında olmadığından serbesttir. $\lambd[x][x x]$
  içinde $x$'in her iki geçişi de $\lambd[x]$ tarafından bağlanır; çünkü
  ikisi de $xx$ içinde serbesttir. $((\lambd[x][xx])x)$ içinde $x$'in son
  geçişi herhangi bir $\lambd[x]$'in kapsamında olmadığından serbesttir.
  $\lambd[x][(\lambd[x][x])x]$ içinde ilk $\lambd[x]$'in kapsamı
  $(\lambd[x][x])x$, ikinci $\lambd[x]$'in kapsamı ise $x$'in sondan bir
  önceki geçişidir. $(\lambd[x][x])x$ içinde $x$'in son geçişi serbest,
  sondan bir önceki geçişi bağlıdır. Dolayısıyla
  $\lambd[x][(\lambd[x][x])x]$ içindeki sondan bir önceki $x$ geçişi ikinci
  $\lambd[x]$, son geçişi ise ilk $\lambd[x]$ tarafından bağlanır.
\end{ex}

Bir $P$ terimindeki bütün değişken geçişlerini inceleyerek serbest değişkenler
kümesini elde edebiliriz. Bu küme $\FV{P}$ ile gösterilir ve doğal biçimde
şöyle tanımlanır:

\begin{defn}[Bir terimin serbest değişkenleri] \ollabel{def:fv}
  Bir terimin \emph{serbest değişkenler} kümesi tümevarımlı olarak şöyle
  tanımlanır:
  \begin{enumerate}
    \item $\FV{x} = \{x\}$ \ollabel{def:fv1}
    \item $\FV{\lambd[x][N]} = \FV{N} \setminus \{x\}$ \ollabel{def:fv2}
    \item $\FV{PQ} = \FV{P} \cup \FV{Q}$ \ollabel{def:fv3}
  \end{enumerate}
\end{defn}

\begin{prob}
  \begin{enumerate}
  \item $\lambd[g][(\lambd[x][g (x x)]) \lambd[x][g (x x)]]$ terimindeki
    $\lambd[g]$ ile iki $\lambd[x]$'in kapsamlarını belirleyin.
  \item $\lambd[g][(\lambd[x][g (x x)]) \lambd[x][g (x x)]]$ içindeki
    bütün değişken geçişleri bağlı mıdır? Her biri hangi soyutlama tarafından
    bağlanır?
    \item $\FV{\lambd[x][(\lambd[y][(\lambd[z][x y]) z]) y]}$ kümesini
    verin.
  \end{enumerate}
\end{prob}

\begin{explain}
Serbest değişken dış dünyaya, yani \emph{ortama}, yapılan bir gönderme
gibidir. Serbest değişkenler içeren bir terim, davranışı ortamın nasıl
kurulduğuna bağlı olduğundan kısmen belirlenmiş bir terim sayılabilir.
Örneğin bir $x$ argümanı alıp bu argümanın $f$ altındaki değerini döndüren
$\lambd[x][f x]$ teriminde $f$ değişkeni serbesttir. Terimin değeri bulunduğu
ortama, özellikle de o ortamdaki $f$ değerine bağlıdır.

Bu terime soyutlama uygularsak $\lambd[f][\lambd[x][f x]]$ elde ederiz. Bu
terim artık ortam değişkeni $f$'ye bağlı değildir; çünkü iki argüman alan ve
birincisini ikinciye uygulamanın sonucunu döndüren bir fonksiyon gösterir.
Ortamda $f$'nin değerini değiştirmek bu terimin davranışını etkilemez: terim,
ortamdaki $f$ değerini değil, yalnızca argüman olarak verilen değeri kullanır.
\end{explain}

\begin{defn}[Kapalı terim, kombinatör]
  Hiçbir serbest değişkeni olmayan terime \emph{kapalı terim} ya da
  \emph{kombinatör} denir.
\end{defn}

\begin{lem}\ollabel{lem:fv}
  \begin{enumerate}
    \item \ollabel{lem:fv-abs} $y \neq x$ ise $y \in
      \FV{\lambd[x][N]}$ ancak ve ancak $y \in \FV{N}$ ise.
    \item \ollabel{lem:fv-app} $y \in \FV{PQ}$ ancak ve ancak
      $y \in \FV{P}$ ya da $y \in \FV{Q}$ ise.
    \end{enumerate}
\end{lem}

\begin{proof}
  Alıştırma.
\end{proof}

\begin{prob}
  \olref[lam][syn][fv]{lem:fv} lemasını ispatlayın.
\end{prob}

